\section{Crossbar Model}\label{sec:model}

We wrote a PyTorch simulator that replaces each linear layer of a model with crossbars and models the
stored value of every cell. This section describes how a weight matrix is stored, the error sources,
and how the simulator was verified. \Cref{tab:model} lists the parameters.

\begin{table}[t]
\centering\footnotesize
\caption{Crossbar model parameters.}\label{tab:model}
\setlength{\tabcolsep}{4pt}
\begin{tabular}{@{}ll@{}}
\toprule
Parameter & Value \\
\midrule
Weight precision & 8-bit signed, one scale per output column \\
Input precision & 8-bit signed \\
Cells per weight (per sign) & 7 SLC ($b=1$) or 4 MLC ($b=2$) \\
Conductance range & $G_{\max}=1$ (normalised), $G_{\max}/G_{\min}=100$ \\
Tile size & $128\times128$ \\
Programming noise $\sigma$ & 0--10\% of $G_{\max}$, fixed per chip \\
Read noise $\sigma_r$ & 0.5--5\% of $G_{\max}$, fresh at every read \\
Stuck cells $p$ & 0.1--5\% of cells, half at $G_{\max}$, half at $G_{\min}$ \\
ADC & ideal, or 3--8 bits with a fixed calibrated range \\
Wire resistance $r$ & $R_{\text{wire}}$ per segment / $R_{\text{on}}$, $10^{-4}$--$3\times10^{-3}$ \\
\bottomrule
\end{tabular}
\end{table}

\subsection{Storing a Weight Matrix}
A matrix $W\in\mathbb R^{m\times n}$, with $n$ inputs and $m$ outputs, is stored with its inputs on the rows
(word lines) and its outputs on the columns (bit lines) of the crossbar. It is quantised to signed 8-bit
integers with one scale per output column,
\begin{equation}
q_{ij}=\operatorname{clip}\!\big(\operatorname{round}(w_{ij}/s_i),\,-127,\,127\big),
\label{eq:quant}
\end{equation}
with column scale $s_i=\max_j|w_{ij}|/127$ for output $i$.
The sign is handled by a differential pair, $q=q^{+}-q^{-}$ with $q^{\pm}=\max(\pm q,0)$, and each
7-bit magnitude is split into $S=\lceil 7/b\rceil$ base-$L$ digits,
\begin{equation}
q^{\pm}=\sum_{t=0}^{S-1}L^{t}d_t^{\pm},\qquad d_t^{\pm}\in\{0,\dots,L-1\},
\label{eq:slice}
\end{equation}
each stored in its own cell at conductance $G=G_{\min}+d\,\Delta$. An SLC weight thus uses $2\times7$
cells and an MLC weight $2\times4$. The stored weight is read back as
\begin{equation}
\hat w_{ij}=s_i\sum_{t=0}^{S-1}L^{t}\,\frac{G^{+}_{ij,t}-G^{-}_{ij,t}}{\Delta},
\label{eq:readback}
\end{equation}
where the $G_{\min}$ offset cancels in the difference. Without errors, $\hat w_{ij}=s_iq_{ij}$ exactly.
Inputs are quantised to 8 bits in the same way with one scale per tensor, and the layer output is the
product of the quantised input with $\hat W$.

\subsection{Error Sources}
Each error source is a separate, seeded option, so that its effect can be measured on its own.
\begin{itemize}
\item \textbf{Programming noise.} Every cell receives a fixed Gaussian offset,
$\tilde G=\max\!\big(0,\;G+\sigma G_{\max}\varepsilon\big)$ with $\varepsilon\sim\mathcal N(0,1)$, drawn
once per chip. Because $\Delta$ is three times smaller for MLC, the same $\sigma$ causes a three times
larger digit error in MLC.

\item \textbf{Level-dependent noise.} The noise of an MLC cell can depend on its level: $\sigma$ is
multiplied by a factor $m_d$ per level $d$. We use $m=(1,1.5,1.5,1)$ and $m=(1,2,2,1)$, making the two
middle levels up to twice as noisy as the end levels.

\item \textbf{Read noise.} An independent Gaussian offset of standard deviation $\sigma_r G_{\max}$ is
added to every cell at every read. Since it is independent of the stored value, its effect on an output
is Gaussian with a closed-form variance, which the simulator adds directly to the output.

\item \textbf{Stuck cells.} A cell is stuck with probability $p$, at $G_{\max}$ or $G_{\min}$ with equal
probability, and keeps that conductance whatever is written. In the \emph{clustered} pattern, faults
grow from random seeds into $3\times3$ blocks, with the seed rate set so that the overall rate is about
$p$. In the \emph{line} pattern, a whole bit line within one $128$-row tile fails with probability $p$.

\item \textbf{ADC and bit-serial inputs.} With a finite ADC, inputs are applied one bit at a time:
each input is written in sign-magnitude form, and bit $k$ of its magnitude is applied with the input's
sign. For every input bit $k$, weight slice $t$ and $128$-row tile, the column sum is quantised by a
$b_{\text{ADC}}$-bit ADC and the results are combined digitally,
\begin{equation}
y_j=\sum_{k}\sum_{t}2^{k}L^{t}\,\mathrm{ADC}_t\!\Big(\sum_{i\in\text{tile}}x_i^{(k)}\big(d^{+}_{ij,t}-d^{-}_{ij,t}\big)\Big).
\label{eq:adc}
\end{equation}
The full-scale range of each slice's ADC is fixed once, from the largest column sum observed on 64
training sentences on an error-free chip, and is not adjusted at run time.

\item \textbf{IR drop.} A first-order static model attenuates each cell by its distance from the row
driver and the column sense amplifier. A cell crossed by $n$ wire segments behaves as
\begin{equation}
G_{\text{eff}}=\frac{G}{1+r\,G\,n},\qquad n=(c+1)+(127-\rho),
\label{eq:irdrop}
\end{equation}
where $c$ and $\rho$ are the cell's column and row positions within its tile, and $r$ is the resistance
of one wire segment relative to the on-resistance. This ignores the dependence of IR drop on the input
pattern, which is enough to show whether a cell's position matters.
\end{itemize}

\subsection{Chips and Paired Comparison}
A \emph{chip} is one draw of all random errors in a model: noise offsets, stuck cells and their values.
All draws come from a single random generator seeded with the chip number, and every placement or
mapping is evaluated on the same seeds. Two placements are therefore compared on chips with identical
error draws, and their difference on each chip is a paired measurement. We report the mean accuracy
over chips and, where differences are small, the mean paired difference with its standard error.

\subsection{Verification}
Every experiment starts with self-tests on random matrices, and stops if any check fails:
\begin{itemize}
\item With all errors disabled, the crossbar output equals the 8-bit integer product to a relative
error below $10^{-4}$, for SLC and MLC, both in one pass and when read slice by slice, tile by tile and
bit by bit.
\item The bit planes of an input rebuild it exactly.
\item A calibrated 16-bit ADC is lossless to a relative error below $10^{-3}$, while a 4-bit ADC is not.
\item The measured spread of read noise matches its closed-form variance to within 15\%.
\item The three fault patterns produce the requested share of stuck cells.
\item IR drop only reduces stored conductances, and level-dependent noise with all factors equal to one
reproduces plain programming noise exactly.
\end{itemize}
End to end, a chip with no errors reaches within about one point of the digital model for every placement
(\cref{sec:results}).
